\section{The Church's Old Testament: Codices, Canon, and the Vulgate Quarrel}

\subsection{The great codices}

The fourth century gave the Greek Bible its monumental form: complete
parchment Bibles produced by Christian scriptoria. Codex Vaticanus and
Codex Sinaiticus (both 4th century) and Codex Alexandrinus (5th) are the
oldest surviving attempts to put the whole of Scripture between two
covers---and they do not agree on what ``the whole'' is. Sinaiticus
includes 1 and 4~Maccabees but not 2--3; Alexandrinus carries all four
books of Maccabees and appends Psalm~151 and a collection of liturgical
Odes; Vaticanus has no Maccabees at all.\endnote{Contents as given in the
standard catalog descriptions of the three codices (see References). The
codices' New Testament supplements (Barnabas and Hermas in Sinaiticus;
the Clementine epistles in Alexandrinus's table) tell the same story of
soft edges.} Two conclusions follow, one positive, one negative. These
pandects witness what fourth- and fifth-century churches actually copied
as Scripture---the translated Hebrew books together with Tobit, Judith,
Wisdom, Sirach, Baruch, Maccabees, and the expanded Esther and Daniel.
They are not an inventory of what Alexandrian Judaism held canonical six
centuries earlier, and projecting their contents backward as ``the
Alexandrian canon'' outruns the evidence.

\subsection{The canon question}

Christian lists tell a story of unresolved tension long before the
Reformation made it urgent. Melito of Sardis (c.~170), who traveled east
to learn accurately the number and order of the books of the old covenant,
returned a list close to the Hebrew Bible---without Esther, and without
Tobit, Judith, Sirach, or Maccabees.\endnote{Melito in Eusebius, \work{Church History} 4.26 (NPNF
series~2), checked at New Advent. The ``Wisdom'' in his list is ambiguous
---possibly an alternative title for Proverbs---and is not pressed
here.} Athanasius's festal letter of 367 canonized twenty-two books
matching the Hebrew count (Esther again excluded; Baruch and the Letter
of Jeremiah counted with Jeremiah) while naming Wisdom, Sirach, Esther,
Judith, and Tobit as books ``appointed by the Fathers to be read'' though
not in the canon.\endnote{Athanasius, \work{Festal Letter} 39 (NPNF
series~2, vol.~4), checked at New Advent.} The African councils of the
390s, with Augustine, listed the fuller church canon; Jerome, in the
same years, planted the flag of the Hebrew count and named the rest
``apocrypha.'' The Catholic resolution came only at Trent: the decree of
8 April 1546 defined the canon with Tobit, Judith, Wisdom,
Ecclesiasticus, Baruch, and both books of Maccabees, ``entire with all
their parts,'' under anathema, and declared the ``old and vulgate
edition'' authentic for public use.\endnote{Council of Trent, Session IV
(8 April 1546), decree on the canonical Scriptures, checked at
papalencyclicals.net. Trent's act settles the Catholic canon question as
law; it does not, and does not claim to, rewrite the earlier history
just narrated. The books are called deuterocanonical in Catholic usage;
Protestant usage retained Jerome's ``apocrypha.''}

\subsection{Jerome against Augustine}

The Septuagint's status in the Latin West was decided in a correspondence.
Jerome, translating the Old Testament afresh from the Hebrew
(c.~390--405), dismantled the origin legend with philological relish: ``I
don't know who was the first author to construct with his lying the
seventy cells in Alexandria,'' he wrote in the preface to his Pentateuch,
when Aristeas and Josephus report the translators gathered in one
basilica, comparing---not prophesying: ``It is one thing to be a seer,
another to be an interpreter''---and he pointed out that the apostles
quote words from the Hebrew that the Seventy do not
contain.\endnote{Jerome, Prologue to Genesis (to Desiderius), checked in
the public-domain Edgecomb translation at tertullian.org against the
NPNF presentation of Jerome's prefaces; \work{Letter} 112 for what
follows, checked at New Advent (= Augustine, \work{Letters} 75).}
Augustine, receiving the new version with alarm, urged Jerome to
translate from the Seventy instead: the Greek Bible was the bond between
Latin and Greek churches, and novelty had costs---at Oea in
Tripolitania, he reported, a congregation rioted when Jonah's gourd
became ivy, the local Jews (consulted as referees) sided against
Jerome's rendering, and the bishop had to retreat.\endnote{Augustine,
\work{Letter} 71 to Jerome, checked at New Advent; his considered
theology of the two texts is \work{City of God} 18.43: one Spirit spoke
through prophets and translators, and divergences may both be prophetic. Jerome's reply defends \textit{hedera} (ivy) for the Hebrew
\textit{ciceion} and denies condemning the Seventy, noting he had
himself earlier revised a Latin Psalter from Origen's hexaplaric Greek,
signs and all.} The outcome was a compromise neither man designed:
Jerome's Hebrew-based version became the Vulgate of the West---carrying
within it his Gallican Psalter, translated not from the Hebrew but from
the hexaplaric Septuagint, and the deuterocanonical books he had
disparaged---while the Greek East simply kept the Seventy. To this day
the Septuagint remains the received Old Testament of the Orthodox
churches, which have never replaced it with a translation from the
Hebrew nor defined a single critical text of it.\endnote{Class~S
synthesis; the Orthodox reception is stated at the level of received
liturgical use, not of any conciliar definition, and modern Orthodox
scholarship engages the critical editions described below.}
